自然言語処理を用いた実用システムの開発にも取り組んできました。対象は、政策決定支援のための将来事象予測、議論マイニングに基づく意思決定支援、エンジニア向けの社内資産 RAG、LLM エージェントによるデジタル回路自動設計などです。

例えば、国家・企業トップの意思決定を支援する将来事象予測技術では、過去の膨大なテキストデータから事象間の因果関係を抽出し、現在事象と因果関係の強い後続事象を提示しました。文書フィルタ、エンティティリンカー、述語項構造解析、事象同一化、因果関係抽出などを組み合わせ、大規模分散処理環境でシステムを構築しました。

また、会議活性化アプリケーションでは、音声認識後の書き起こしテキストを基に、話者分類、意図推定、スレッド分割などをリアルタイムで行い、ユーザに提示しました。このほかにも、プロジェクト固有の文脈を踏まえて有用な過去資産を検索する社内資産 RAG、企業やプロジェクト固有の文脈を踏まえたコード検索システム CHICOT、LLM エージェントによるデジタル回路自動設計など、複数の実用システムを開発してきました。
